\section{Introduction}\label{sec:introduction}

In a broadcast process on a tree, each vertex transmits a noisy copy of
its spin to its children. The reconstruction problem asks whether the
spins far from the root retain information about its initial spin.
For the symmetric four-state channel, the linear second-moment threshold
has long been expected to be the exact reconstruction threshold in the
ferromagnetic regime. We prove the nonreconstruction direction at and
below this threshold for all regular and Poisson branching parameters.

\subsection{Model and main result}
Write $[4]=\{1,2,3,4\}$. The root spin $\sigma_\rho$ is uniform on $[4]$.
For a parameter $\lambda\in[0,1]$, each edge independently transmits the
parent spin through the channel
\begin{equation}\label{eq:broadcast-channel}
 P_\lambda(j\mid i)
   =\lambda\mathbf 1_{\{i=j\}}+\frac{1-\lambda}{4},
 \qquad i,j\in[4].
\end{equation}
The parameter $\lambda$ is the nontrivial eigenvalue of the channel;
its nonnegativity specifies the ferromagnetic regime.

We consider the rooted $d$-ary tree, where every vertex has exactly $d$
children, and a Galton--Watson tree with independent Poisson$(d)$ offspring
counts. In the latter case the tree is sampled independently of the root
spin and the randomness used by the edge channels. The
observer is given the entire rooted tree $T$ and the spins
$\sigma_{L_n}$ at level $L_n=\{v:\operatorname{dist}(v,\rho)=n\}$.
Let
\begin{equation}\label{eq:advantage}
 a_n(d,\lambda)=
 \E\left[\frac12\sum_{i=1}^4
 \left|\Pr(\sigma_\rho=i\mid T,\sigma_{L_n})-\frac14\right|\right].
\end{equation}
The tree observation in the Galton--Watson model is unlabelled, the level
observation is empty when $L_n$ is empty, and the expectation is not
conditioned on survival. Nonreconstruction means that $a_n(d,\lambda)$
tends to zero.

\begin{theorem}\label{thm:main}
For the channel \eqref{eq:broadcast-channel}, if $d\lambda^2\leq1$, then
\[
 \lim_{n\to\infty}a_n(d,\lambda)=0
\]
in each of the following models:
\begin{enumerate}
 \item the rooted $d$-ary tree, for every integer $d\geq2$;
 \item the Poisson$(d)$ Galton--Watson tree, for every real $d>0$.
\end{enumerate}
In particular, the assertion includes $d\lambda^2=1$.
\end{theorem}

Here $d$ denotes the number of children or its mean, not the total
undirected degree. The known Kesten--Stigum reconstruction direction for
$d\lambda^2>1$ makes Theorem~\ref{thm:main} an exact threshold statement
\cite{KestenStigum1966,MosselPeres2003,Lyons1990}. For Poisson trees,
the branching-number reconstruction criterion applies on survival,
where the branching number equals $d$; survival has positive probability
when $d\lambda^2>1$.
The theorem is computer assisted: all probabilistic arguments are proved
below, and the remaining polynomial inequalities are verified by the
complete exact-arithmetic programs in Appendix~\ref{app:verifiers}.

\subsection{History and significance}
The Kesten--Stigum criterion originates in the second-moment analysis of
multitype branching processes \cite{KestenStigum1966}. In tree broadcasting it gives the
reconstruction direction above $d\lambda^2=1$, but the corresponding
nonreconstruction assertion depends on the channel. For binary symmetric
channels, Bleher, Ruiz, and Zagrebnov established sharpness on regular
trees \cite{BleherRuizZagrebnov1995}; Evans, Kenyon, Peres, and Schulman
developed the general-tree theory
\cite{EvansKenyonPeresSchulman2000}. Mossel showed that the spectral
criterion need not be necessary for other channels \cite{Mossel2001}.
For Potts channels, this distinction becomes decisive as the number of
states increases. M\'ezard and Montanari proposed threshold coincidence
for small alphabets and discussed its extension to all regular degrees
\cite[Section~8]{MezardMontanari2006}.
Sly proved sharpness for the three-state model on regular trees of
sufficiently large degree and established failure of sharpness for
$q\geq5$ \cite{Sly2011}. The four-state case lies at the boundary between
these behaviors.

Mossel, Sly, and Sohn proved nonreconstruction, including the critical
point, for the three- and four-state models on Galton--Watson trees of
sufficiently large mean degree under hypotheses that cover both regular
and Poisson offspring \cite{MosselSlySohn2025}.
Theorem~\ref{thm:main} removes the large-degree restriction for the
ferromagnetic four-state model in these two families. Cavity expansions
and numerical investigations by Ricci-Tersenghi, Semerjian, and
Zdeborov\'a also support the four-state ferromagnetic threshold
\cite{RicciTersenghiSemerjianZdeborova2019}; those predictions are distinct
from a rigorous all-degree nonreconstruction theorem. The restriction
$\lambda\geq0$ matters: the four-state antiferromagnetic model can
reconstruct below the Kesten--Stigum threshold
\cite{MosselSlySohn2025}.

\subsection{The proof mechanism}
The key object is the distribution of a posterior probability vector
$p=(p_1,p_2,p_3,p_4)$, not an individual vector. Define its quadratic
information by
\[
 A(p)=4\sum_{i=1}^4p_i^2-1.
\]
A channel step multiplies $\E A$ by $\lambda^2$. The problem is to control
the nonlinear combination of the observations arriving from different
children. We construct a degree-six polynomial $H$ for which the
additional constraint $\E H(p)\geq0$ is preserved by both channel steps
and the combination of observations. Two explicit polynomial inequalities
then make $\E A$ subadditive under combination, with a strictly positive
cubic saving when two informative branches are present.

Writing $m_n$ for the expected quadratic information at depth $n$, this
gives, for $\lambda>0$, a recursion of the form
\[
 m_{n+1}\leq d\lambda^2m_n-cm_n^3,
 \qquad c=c(d,\lambda)>0.
\]
The cubic term is essential at equality: it forces $m_n\to0$ even when
the linear coefficient equals one. The constraint on the posterior law
is linear, so it also survives mixtures over the observed offspring
count. This is what allows the same pairwise inequalities to treat the
Poisson and regular models.

Expected functionals of posterior distributions are central to earlier
entropy and stochastic Lyapunov methods, including the work of Formentin
and K\"ulske \cite{FormentinKulske2009}. Rigorous computational methods
for nonreconstruction include Bhatnagar and Maneva's survey method
\cite{BhatnagarManeva2011} and Gu's population-dynamics method for binary
Poisson hypertrees \cite{Gu2026}. The additional ingredient here is the
explicit polynomial constraint that remains valid during each pairwise
combination of observations. Its pairing with the two combination
inequalities gives a strict moment saving for every parameter in the two
branching families when $\lambda>0$.

The required inequalities are certified on products of
probability simplices by nonnegative homogeneous coefficients, using
Taylor lower bounds and interpolation with explicitly certified
remainders. The interpolation itself is never assumed to approximate
the target accurately. Thus the computation verifies fixed polynomial
inequalities on all posterior vectors; the preserved constraint on their
laws then supplies the tree argument.

Section~\ref{sec:experiments} gives the posterior operations, and
Section~\ref{sec:cone} states the explicit polynomial inequalities and
derives their consequences for posterior laws.
Section~\ref{sec:tree-decay} completes the probabilistic argument.
Sections~\ref{sec:polynomial-certificates} and
\ref{sec:finite-verification} prove the inequalities through a finite
exact-arithmetic calculation.
